\documentclass{article}
\usepackage{amsmath,amssymb,amsthm}
\usepackage{algorithm}
\usepackage{algpseudocode}
\usepackage{bm}
\usepackage{hyperref}
\newtheorem{theorem}{Theorem}
\newtheorem{lemma}{Lemma}
\newtheorem{assumption}{Assumption}
\newtheorem{corollary}{Corollary}
\begin{document}

% -------------------------------------------------
\section*{Heavy‑Ball Momentum (HBM)}
% -------------------------------------------------

\section*{Mathematical Formulation}
Let $f:\mathbb{R}^d\to\mathbb{R}$ be a differentiable objective function.
The Heavy‑Ball method with constant stepsize $\alpha>0$ and momentum
parameter $\beta\in[0,1)$ generates a sequence $\{x_k\}_{k\ge 0}$
according to  

\[
\boxed{
\begin{aligned}
x_{k+1} &= x_k + \beta\bigl(x_k-x_{k-1}\bigr)-\alpha \nabla f(x_k),
\qquad k\ge 1,\\[2mm]
x_1 &= x_0 - \alpha \nabla f(x_0) \quad\text{(first step, no momentum)} .
\end{aligned}}
\]

The hyper‑parameters are  

\begin{itemize}
  \item stepsize (learning rate) $\alpha>0$,
  \item momentum coefficient $\beta\in[0,1)$.
\end{itemize}

\section*{Pseudocode}
\begin{algorithm}[H]
\caption{Heavy‑Ball Momentum}
\begin{algorithmic}[1]
\Require Initial point $x_0\in\mathbb{R}^d$, stepsize $\alpha>0$, momentum $\beta\in[0,1)$,
        maximum iterations $K$.
\State $x_{-1}\gets x_0$ \Comment{Define $x_{-1}$ for uniform notation}
\State Compute $g_0 \gets \nabla f(x_0)$
\State $x_1 \gets x_0 - \alpha g_0$ \Comment{First iteration, no momentum}
\For{$k=1$ \textbf{to} $K-1$}
    \State $g_k \gets \nabla f(x_k)$
    \State $x_{k+1} \gets x_k + \beta (x_k-x_{k-1}) - \alpha g_k$
\EndFor
\State \Return $x_K$
\end{algorithmic}
\end{algorithm}

\section*{Dry Run Example}
Consider the univariate quadratic  

\[
f(w) = (w-3)^2 .
\]

Its gradient is $\nabla f(w)=2(w-3)$.  
Choose $x_0=0$, stepsize $\alpha=0.1$, momentum $\beta=0.5$.
We compute three updates.

\begin{center}
\begin{tabular}{c|c|c|c}
$k$ & $x_k$ & $\nabla f(x_k)$ & $f(x_k)$ \\ \hline
0 & $0$ & $2(0-3)=-6$ & $9$ \\[1mm]
1 & $x_1 = 0 -0.1(-6)=0.6$ & $2(0.6-3)=-4.8$ & $5.76$ \\[1mm]
2 & $x_2 = 0.6 +0.5(0.6-0) -0.1(-4.8)=0.6+0.3+0.48=1.38$ &
$2(1.38-3)=-3.24$ & $2.6244$ \\[1mm]
3 & $x_3 = 1.38 +0.5(1.38-0.6) -0.1(-3.24)=1.38+0.39+0.324=2.094$ &
$2(2.094-3)=-1.812$ & $0.820$ \\ 
\end{tabular}
\end{center}

We see the iterates rapidly approach the minimizer $w^\star=3$ and the
objective value decreases.

% -------------------------------------------------
\section*{Assumptions}
% -------------------------------------------------
\begin{assumption}[Strong Convexity]\label{as:sc}
$f$ is $\mu$‑strongly convex, i.e. for all $x,y\in\mathbb{R}^d$  

\[
f(y)\ge f(x)+\langle \nabla f(x),y-x\rangle+\frac{\mu}{2}\|y-x\|^2,
\qquad\mu>0 .
\]
\end{assumption}

\begin{assumption}[Smoothness]\label{as:smooth}
$f$ has $L$‑Lipschitz continuous gradient, i.e.  

\[
\|\nabla f(x)-\nabla f(y)\|\le L\|x-y\|,\qquad L>0 .
\]
\end{assumption}

\begin{assumption}[Parameter Range]\label{as:param}
The stepsize $\alpha$ and momentum $\beta$ satisfy  

\[
0<\alpha<\frac{2(1+\beta)}{L},\qquad
0\le\beta<1 .
\]
\end{assumption}

All results below are derived under Assumptions~\ref{as:sc}–\ref{as:param}.

% -------------------------------------------------
\section*{Main Convergence Theorem}
% -------------------------------------------------
\begin{theorem}[Linear convergence of Heavy‑Ball]\label{thm:hb}
Let $\{x_k\}$ be generated by the Heavy‑Ball iteration.
Define the error vector $e_k:=x_k-x^\star$, where $x^\star$ is the unique
minimizer of $f$.
If the parameters $\alpha,\beta$ satisfy  

\[
0<\alpha<\frac{2}{L},\qquad
0\le \beta<\frac{(1-\sqrt{\alpha\mu})^2}{1+\sqrt{\alpha\mu}},
\tag{*}
\]

then there exists a constant $\rho\in(0,1)$ such that  

\[
\|e_{k}\| \le \rho^{k}\,\|e_{0}\|,\qquad \forall k\ge 0 .
\]

Equivalently, the function values converge linearly:

\[
f(x_k)-f(x^\star) \le \frac{L}{2}\rho^{2k}\|e_{0}\|^{2}.
\]
\end{theorem}

\section*{Theoretical Properties}
\begin{itemize}
  \item \textbf{Stability.} Condition~$(*)$ guarantees that the spectral radius
        of the iteration matrix is $<1$, which is necessary and sufficient
        for asymptotic stability of the linearized dynamics.
  \item \textbf{Hyper‑parameter sensitivity.}
        The admissible region for $\beta$ shrinks as $\alpha$ approaches its
        upper bound $2/L$. The optimal choice (in the sense of minimizing $\rho$)
        is  

        \[
        \alpha^\star=\frac{4}{(\sqrt{L}+\sqrt{\mu})^{2}},\qquad
        \beta^\star=\left(\frac{\sqrt{L}-\sqrt{\mu}}{\sqrt{L}+\sqrt{\mu}}\right)^{2},
        \]
        which yields $\rho = \frac{\sqrt{L}-\sqrt{\mu}}{\sqrt{L}+\sqrt{\mu}}$.
  \item \textbf{Comparison to Gradient Descent.}
        Gradient descent with stepsize $2/(L+\mu)$ attains the same linear
        factor $\rho_{\text{GD}} = \frac{L-\mu}{L+\mu}$, while the optimal
        Heavy‑Ball factor $\rho^\star$ is strictly smaller:
        $\rho^\star = \frac{\sqrt{L}-\sqrt{\mu}}{\sqrt{L}+\sqrt{\mu}} < 
        \frac{L-\mu}{L+\mu}$ for $L>\mu$.
\end{itemize}

\section*{Discussion}
The Heavy‑Ball method can be interpreted as a discretization of the
continuous‐time second‑order differential equation
$\ddot{x}(t)+\gamma\dot{x}(t)+\nabla f(x(t))=0$.
The momentum term $\beta(x_k-x_{k-1})$ mimics the velocity,
and the strong convexity–smoothness pair $(\mu,L)$ determines a
stability region identical to that of a damped linear oscillator.
The theorem above shows that, under the classical Polyak
parameter choice, the method enjoys a provably faster linear rate
than plain gradient descent while retaining a simple first‑order
implementation.

% -------------------------------------------------
\section*{Preliminaries (Definitions and Lemmas)}
% -------------------------------------------------
\begin{lemma}[Quadratic upper–lower bounds]\label{lem:quad}
For any $\mu$‑strongly convex and $L$‑smooth $f$ and any $x$,
\[
\frac{\mu}{2}\|x-x^\star\|^{2}\le f(x)-f(x^\star)\le
\frac{L}{2}\|x-x^\star\|^{2}.
\]
\end{lemma}
\begin{proof}
The lower bound follows from the definition of strong convexity with
$y=x^\star$; the upper bound follows from $L$‑smoothness applied to the
gradient at $x^\star$ (where $\nabla f(x^\star)=0$). ∎
\end{proof}

\begin{lemma}[Descent Lemma]\label{lem:descent}
If $\nabla f$ is $L$‑Lipschitz, then for any $x$ and any $d$,
\[
f(x+d)\le f(x)+\langle\nabla f(x),d\rangle+\frac{L}{2}\|d\|^{2}.
\]
\end{lemma}
\begin{proof}
Standard consequence of the integral form of the gradient
and the Lipschitz bound. ∎
\end{proof}

\section*{Main Proof (Step‑by‑Step Derivation)}
We prove Theorem~\ref{thm:hb}.  Throughout the proof we denote
$e_k:=x_k-x^\star$ and use the optimality condition
$\nabla f(x^\star)=0$.

\subsection*{Step 1 – Rewrite the iteration in error form}
\[
\begin{aligned}
e_{k+1}
&=x_{k+1}-x^\star \\
&=x_k-x^\star +\beta\bigl(x_k-x_{k-1}\bigr)-\alpha\nabla f(x_k)\\
&=e_k+\beta(e_k-e_{k-1})-\alpha\nabla f(x_k).
\end{aligned}
\tag{1}
\]

\subsection*{Step 2 – Linearisation using smoothness}
Because $f$ is $L$‑smooth and $\mu$‑strongly convex,
its gradient is a monotone operator satisfying  

\[
\mu I \preceq \nabla^{2}f(\xi) \preceq L I
\quad\text{for some }\xi\text{ on the line segment }[x^\star,x_k].
\]

Hence there exists a symmetric matrix $H_k$ with eigenvalues in $[\mu,L]$ such that  

\[
\nabla f(x_k)=H_k e_k .
\tag{2}
\]

\subsection*{Step 3 – Substitute (2) into (1)}
\[
e_{k+1}= (1+\beta) e_k - \beta e_{k-1} -\alpha H_k e_k
      = \underbrace{\bigl[(1+\beta)I-\alpha H_k\bigr]}_{A_k} e_k
        -\beta e_{k-1}.
\tag{3}
\]

\subsection*{Step 4 – Form a 2‑dimensional state vector}
Define $z_k:=\begin{bmatrix}e_k\\ e_{k-1}\end{bmatrix}\in\mathbb{R}^{2d}$.
Then (3) can be written as  

\[
z_{k+1}=M_k z_k,\qquad 
M_k:=\begin{bmatrix}
A_k & -\beta I\\
I   & 0
\end{bmatrix}.
\tag{4}
\]

\subsection*{Step 5 – Spectral bound on $M_k$}
Since $H_k$ shares the same eigenvectors as $I$, each eigen‑direction
decouples.  Let $\lambda\in[\mu,L]$ be an arbitrary eigenvalue of $H_k$.
In that direction the scalar matrix reads  

\[
\tilde M(\lambda)=
\begin{bmatrix}
1+\beta-\alpha\lambda & -\beta\\
1                    & 0
\end{bmatrix}.
\tag{5}
\]

The eigenvalues $\xi$ of $\tilde M(\lambda)$ satisfy  

\[
\xi^{2}-(1+\beta-\alpha\lambda)\xi+\beta=0.
\tag{6}
\]

\subsection*{Step 6 – Choose $\alpha,\beta$ to make the spectral radius $<1$}
The quadratic (6) has roots inside the unit disc iff  

\[
|\,\beta\,|<1,\qquad
|\,1+\beta-\alpha\lambda\,|<1+\beta .
\tag{7}
\]

The first inequality holds because $0\le\beta<1$ (Assumption~\ref{as:param}).
For the second inequality, the worst case occurs at the extreme
eigenvalues $\lambda\in\{\mu,L\}$.  Enforcing (7) for both yields  

\[
-1-\beta < 1+\beta-\alpha\mu < 1+\beta,
\qquad
-1-\beta < 1+\beta-\alpha L   < 1+\beta .
\]

These simplify to  

\[
0<\alpha<\frac{2(1+\beta)}{L},
\qquad
\alpha > \frac{2\beta}{L+\mu}.
\tag{8}
\]

A convenient sufficient condition (used in many textbooks) is  

\[
0<\alpha<\frac{2}{L},
\qquad
0\le\beta<\frac{(1-\sqrt{\alpha\mu})^{2}}{1+\sqrt{\alpha\mu}} .
\tag{9}
\]

Under (9) the spectral radius of $\tilde M(\lambda)$ is bounded by a
common constant $\rho\in(0,1)$ that does not depend on $k$ (the bound is
obtained by maximizing the modulus of the roots of (6) over
$\lambda\in[\mu,L]$).  Consequently  

\[
\|z_{k+1}\|\le\rho\|z_k\|\quad\Longrightarrow\quad
\|e_k\|\le\rho^{k}\|e_0\|.
\tag{10}
\]

\subsection*{Step 7 – From error norm to function value}
Using Lemma~\ref{lem:quad},  

\[
f(x_k)-f(x^\star)\le \frac{L}{2}\|e_k\|^{2}
                 \le \frac{L}{2}\rho^{2k}\|e_0\|^{2}.
\tag{11}
\]

\subsection*{Conclusion}
Equations (10)–(11) establish the linear convergence claim of
Theorem~\ref{thm:hb}.  ∎

\section*{Rate Derivation (With Constants)}
From the optimal parameters $\alpha^\star,\beta^\star$ given in
Section~\ref{sec:theoretical-properties}, the characteristic polynomial
(6) becomes  

\[
\xi^{2}-\frac{2\sqrt{L\mu}}{L+\mu}\,\xi+\left(\frac{\sqrt{L}-\sqrt{\mu}}
{\sqrt{L}+\sqrt{\mu}}\right)^{2}=0,
\]

whose roots are both equal to  

\[
\rho^\star=\frac{\sqrt{L}-\sqrt{\mu}}{\sqrt{L}+\sqrt{\mu}}.
\]

Thus the optimal Heavy‑Ball rate is  

\[
\|x_k-x^\star\|\le (\rho^\star)^k\|x_0-x^\star\|,
\qquad
f(x_k)-f(x^\star)\le \frac{L}{2}(\rho^\star)^{2k}\|x_0-x^\star\|^{2}.
\]

\section*{Special Cases}
\begin{itemize}
  \item \textbf{Quadratic functions.} When $f(x)=\frac12 x^\top Q x -b^\top x$,
        with $Q$ SPD and eigenvalues in $[\mu,L]$, the matrix $H_k\equiv Q$
        is constant.  The analysis above becomes exact (no inequalities),
        and the bound $\rho$ coincides with the spectral radius of the
        iteration matrix.
  \item \textbf{Zero momentum ($\beta=0$).} The method reduces to plain
        gradient descent with stepsize $\alpha$, and the convergence factor
        becomes $|1-\alpha\mu|$, which is optimal for $\alpha=2/(L+\mu)$.
  \item \textbf{Over‑damped regime ($\beta$ large).} If $\beta$ exceeds the
        bound in (9), one of the roots of (6) leaves the unit disc and the
        iterates diverge.
\end{itemize}

\end{document}